\documentclass[11pt]{article}

\usepackage[margin=1in]{geometry}
\usepackage{amsmath, amssymb}
\usepackage{xcolor}
\usepackage{listings}
\usepackage{hyperref}
\usepackage{array}
\usepackage{booktabs}

\hypersetup{
    colorlinks=true,
    linkcolor=blue!60!black,
    urlcolor=blue!60!black,
    citecolor=blue!60!black
}

\lstdefinestyle{bugpython}{
  language=Python,
  basicstyle=\ttfamily\small,
  keywordstyle=\color{blue!60!black},
  stringstyle=\color{green!40!black},
  commentstyle=\color{gray!70!black},
  showstringspaces=false,
  breaklines=true,
  frame=single,
  rulecolor=\color{gray!30},
  columns=fullflexible,
  keepspaces=true
}

\title{SymPy Bug Report}
\author{}
\date{}

\begin{document}

\maketitle

\section{Bug 28: solveset returns -1 for a principal cube-root equation}

\subsection{Minimal Reproducer}

\medskip

\begin{lstlisting}[style=bugpython]
from sympy import Eq, N, Rational, S, simplify, solveset, symbols

x = symbols("x")
eq = Eq(x**Rational(1, 3), -1)
sol = solveset(eq, x, S.Reals)

print("equation:", eq)
print("solveset:", sol)
for s in sol:
    print("candidate:", s)
    print("lhs at candidate:", eq.lhs.subs(x, s))
    print("residual:", simplify(eq.lhs.subs(x, s) - eq.rhs))
    print("numeric residual:", N(eq.lhs.subs(x, s) - eq.rhs, 30))
print("complex-domain solveset:", solveset(eq, x, S.Complexes))
\end{lstlisting}

\subsection{Observed Behavior}

\medskip

\begin{lstlisting}[style=bugpython]
equation: Eq(x**(1/3), -1)
solveset: {-1}
candidate: -1
lhs at candidate: (-1)**(1/3)
residual: 1 + (-1)**(1/3)
numeric residual: 1.5 + 0.866025403784438646763723170753*I
complex-domain solveset: EmptySet
\end{lstlisting}

SymPy returns the finite set \(\{-1\}\) for the real-domain call
\[
\operatorname{solveset}(\operatorname{Eq}(x^{1/3},-1),x,\mathbb{R}).
\]
However, substituting the returned value into the original SymPy expression
does not satisfy the equation.  The residual at \(x=-1\) is
\[
(-1)^{1/3}+1,
\]
which SymPy numerically evaluates as \(1.5+0.8660254037844386\,i\), not zero.

\subsection{Expected Behavior}

The correct result is \texttt{EmptySet}.  SymPy's expression
\texttt{x**Rational(1, 3)} is a principal complex power, not the real-valued
cube-root function.  Therefore no real \(x\) satisfies the principal-power
equation \(x^{1/3}=-1\) in SymPy's semantics.

\subsection{Mathematical Explanation}

For SymPy's principal complex power,
\[
z^{1/3}=\exp\!\left(\frac{\operatorname{Log}(z)}{3}\right),
\]
where \(\operatorname{Log}\) is the principal logarithm.  Since
\(\operatorname{Log}(-1)=i\pi\), the principal cube root of \(-1\) is
\[
(-1)^{1/3}=\exp(i\pi/3)=\frac{1}{2}+\frac{\sqrt{3}}{2}i.
\]
Thus
\[
(-1)^{1/3}\ne -1
\quad\text{and}\quad
(-1)^{1/3}+1\ne0.
\]

The returned value is not merely an alternate branch representation of the same
solution.  It is the value produced by the real cube-root function
\(\operatorname{real\_root}(-1,3)=-1\), but the input equation was written using
SymPy's principal \texttt{Pow}.  Solving the principal-power equation therefore
requires respecting the principal branch, and \(-1\) is an extraneous solution.

\subsection{Source-Level Diagnosis}

The diagnosis identifies the responsible behavior in
\nolinkurl{sympy/solvers/solveset.py}, in the function
\texttt{\_invert\_real}, around line 284 in the inspected SymPy 1.14.0 source.
The public call enters \texttt{solveset}; because the requested domain is real,
the symbol is recast to a real dummy and the equation is passed to
\texttt{\_solveset} with checking enabled.  \texttt{\_solveset} rewrites the
equation as \texttt{x**(1/3) + 1} and calls \texttt{\_invert}, which dispatches
to \texttt{\_invert\_real} for the real-domain inversion of the rational power.

The key rule identified in \texttt{\_invert\_real} is the odd-denominator
rational-power branch:

\begin{lstlisting}[style=bugpython]
if den % 2 == 1:
    root = Lambda(n, real_root(n, expo))
    res = imageset(root, g_ys)
    ...
    return _invert_real(base, res, symbol)
\end{lstlisting}

For \(x^{1/3}+1=0\), this branch applies \texttt{real\_root} to the target set
\(\{-1\}\), producing \(\{-1\}\) as the preimage for the base \(x\).  That
transformation is valid for equations explicitly expressed with real-root
semantics, but it is not equivalent to inverting SymPy's principal complex
\texttt{Pow} on negative real inputs.

The diagnosis also explains why the later checking pass does not remove the
extraneous candidate.  The finite-solution filtering path uses
\texttt{domain\_check}, which checks whether the expression is finite and
defined at the candidate, rather than re-evaluating whether the original
equation is true.  Since \((-1)^{1/3}+1\) is finite, the candidate survives even
though its residual is nonzero.  A plausible fix direction is to restrict this
odd-denominator \texttt{real\_root} inversion branch to expressions whose
semantics are actually real-root semantics, or to add a principal-branch guard
and equation-truth validation when branch-sensitive inversions have been used.

\subsection{Additional Failing Instances}

The same failure appears for odd denominators beyond \(3\).  The independent
verification cases consider
\[
x^{1/n}=-1
\]
over the real domain for \(n\in\{3,5,7,9,11,13\}\).  In each case, SymPy's
real-domain solver returns \(\{-1\}\), while the expected answer is
\texttt{EmptySet}.  The reason is uniform: for every listed odd \(n>1\), the
principal value is
\[
(-1)^{1/n}=\exp(i\pi/n),
\]
which is not \(-1\).  Hence \(x=-1\) does not satisfy any of the listed
principal-power equations.

\begin{center}
\small
\renewcommand{\arraystretch}{1.3}
\begin{tabular}{@{}>{\raggedright\arraybackslash}p{0.44\linewidth}>{\raggedright\arraybackslash}p{0.23\linewidth}>{\raggedright\arraybackslash}p{0.23\linewidth}@{}}
\toprule
\textbf{Input / Expression} & \textbf{SymPy behavior} & \textbf{Expected behavior} \\
\midrule
\(\operatorname{solveset}(\operatorname{Eq}(x^{1/3},-1),x,\mathbb{R})\) & returns \(\{-1\}\) & \texttt{EmptySet} \\
\(\operatorname{solveset}(\operatorname{Eq}(x^{1/5},-1),x,\mathbb{R})\) & returns \(\{-1\}\) & \texttt{EmptySet} \\
\(\operatorname{solveset}(\operatorname{Eq}(x^{1/7},-1),x,\mathbb{R})\) & returns \(\{-1\}\) & \texttt{EmptySet} \\
\(\operatorname{solveset}(\operatorname{Eq}(x^{1/9},-1),x,\mathbb{R})\) & returns \(\{-1\}\) & \texttt{EmptySet} \\
\(\operatorname{solveset}(\operatorname{Eq}(x^{1/11},-1),x,\mathbb{R})\) & returns \(\{-1\}\) & \texttt{EmptySet} \\
\(\operatorname{solveset}(\operatorname{Eq}(x^{1/13},-1),x,\mathbb{R})\) & returns \(\{-1\}\) & \texttt{EmptySet} \\
\bottomrule
\end{tabular}
\end{center}

\subsection{Related Issues Assessment}

The bug-finding harness performed a best-effort deduplication check and
classified this as a new instance of a known failure mode.  The closest public
materials identified were SymPy documentation describing the explicit domain
argument to \texttt{solveset}, principal-branch behavior for elementary
functions, and the distinction between principal powers and real-root
interpretations.  The available evidence did not identify a public GitHub
issue, pull request, Stack Overflow question, mailing-list post, or release-note
entry for this exact reproducer returning \(\{-1\}\).  Although rational powers
and principal branches are a known source of solver mistakes, this concrete
real-domain equation is previously unreported in the available evidence and
remains individually actionable as an issue and as a regression test.

\subsection{Proposed SymPy Regression Test}

\medskip

\begin{lstlisting}[style=bugpython]
import pytest

from sympy import Eq, EmptySet, Rational, S, solveset, symbols


@pytest.mark.parametrize("denominator", [3, 5, 7, 9, 11, 13])
def test_principal_odd_root_equation_has_no_negative_real_solution(denominator):
    x = symbols("x")
    assert solveset(Eq(x**Rational(1, denominator), -1), x, S.Reals) == EmptySet
\end{lstlisting}

\end{document}
